\chapter{De Quervain Release}

\section*{What Is This Surgery?}
De Quervain release is a focused surgical procedure designed to alleviate the painful symptoms of De Quervain tenosynovitis, a common inflammatory condition affecting the tendons on the thumb side of the wrist. This condition arises from the constriction and inflammation of the abductor pollicis longus (APL) and extensor pollicis brevis (EPB) tendons as they pass through a tight fibro-osseous tunnel, known as the first dorsal compartment, at the radial styloid. The surgical intervention involves precisely dividing the roof of this compartment, thereby decompressing the entrapped tendons and allowing them to glide freely without friction or impingement. It is typically reserved for patients whose symptoms are persistent and debilitating despite a comprehensive trial of nonoperative management, including splinting, corticosteroid injections, and activity modification. Successful release aims to provide lasting pain relief, restore full and comfortable thumb and wrist function, and prevent recurrence, making it a highly effective and frequently performed hand surgery for this specific pathology.

\section*{Alternative Names}
\begin{itemize}
  \item First dorsal compartment release.
  \item De Quervain's tenosynovitis surgery.
  \item Release of the first extensor compartment.
  \item Decompression of the abductor pollicis longus and extensor pollicis brevis tendons.
  \item Radial styloid tenosynovectomy (when inflamed synovium is also excised).
  \item Extensor pollicis brevis and abductor pollicis longus release.
\end{itemize}

\section*{Relevant Surgical Anatomy}
The first dorsal compartment of the wrist is a critical anatomical structure located on the radial aspect of the distal forearm, overlying the radial styloid. This fibro-osseous tunnel is formed by the radial styloid bone dorsally and the overlying extensor retinaculum, a thickened fibrous band, volarly. Within this compartment, two distinct tendons are housed: the abductor pollicis longus (APL) and the extensor pollicis brevis (EPB). The APL tendon, typically larger, originates from the posterior surfaces of the ulna, radius, and interosseous membrane, inserting into the base of the first metacarpal. The EPB tendon, often smaller and sometimes absent, originates from the posterior surface of the radius and interosseous membrane, inserting into the base of the proximal phalanx of the thumb. These tendons are encased in a common synovial sheath, which can become inflamed and thickened in tenosynovitis.

Crucially, the superficial radial sensory nerve (SRSN) and its branches are in close proximity to the first dorsal compartment, typically coursing radially and dorsally over the retinaculum. These nerve branches provide sensation to the dorsum of the thumb, index, middle, and radial half of the ring fingers, and are highly susceptible to injury during surgical dissection. The radial artery, while deeper and more volar, is also a structure to be mindful of, though less directly at risk than the SRSN. Surgical landmarks include the palpable radial styloid and the anatomical snuffbox, which is bordered by the APL and EPB tendons radially and the extensor pollicis longus (EPL) ulnarly. Variations in anatomy, such as the presence of a separate subcompartment for the EPB tendon or accessory slips of the APL, are common and must be identified during surgery to ensure complete decompression.

\section*{Surgical Principle \& Rationale}
The fundamental surgical principle of De Quervain release is the mechanical decompression of the entrapped abductor pollicis longus (APL) and extensor pollicis brevis (EPB) tendons within the first dorsal compartment. The underlying pathology involves a mismatch between the volume of the tendons and their synovial sheath, and the constricted volume of the fibro-osseous tunnel, exacerbated by inflammation and thickening of the retinacular roof. Repetitive thumb and wrist movements lead to friction, microtrauma, and chronic inflammation, which further narrows the compartment and compresses the tendons, resulting in pain and impaired function.

The rationale for surgical intervention is to permanently enlarge this constricting tunnel by longitudinally incising the thickened extensor retinaculum that forms its roof. This division of the retinaculum immediately relieves the pressure on the APL and EPB tendons, allowing them to glide freely and smoothly without impingement. Technical goals include ensuring complete release of all tendon slips and any accessory compartments, meticulously identifying and protecting the superficial radial sensory nerve branches to prevent iatrogenic injury, and avoiding excessive release that could lead to tendon subluxation or bowstringing. In some cases, inflamed synovial tissue may also be excised (synovectomy) to reduce inflammation and prevent recurrence, although the primary goal remains mechanical decompression.

\section*{History \& Surgical Evolution}
De Quervain tenosynovitis was first comprehensively described in 1895 by the Swiss surgeon Fritz de Quervain, who meticulously documented a series of patients presenting with characteristic painful swelling and tenderness on the radial side of the wrist. Prior to this, similar conditions may have been recognized but lacked a specific anatomical understanding or nomenclature. Early treatments were largely conservative, focusing on rest, immobilization, and topical remedies, often with limited success for chronic cases.

As the anatomical understanding of the first dorsal compartment and the pathophysiology of tendon entrapment evolved in the early 20th century, surgical approaches began to emerge. The concept of surgically dividing the constricting retinaculum to decompress the tendons gained traction. By the mid-20th century, surgical release had become the established definitive treatment for patients with refractory De Quervain tenosynovitis. Subsequent decades saw refinements in surgical technique, particularly concerning incision placement (moving towards smaller, more cosmetic transverse incisions), meticulous dissection to identify and protect the superficial radial sensory nerve, and careful inspection for anatomical variations such as accessory compartments. The core principle of retinacular division has remained constant, but improved understanding of nerve anatomy and variations has significantly reduced complications and enhanced the safety and efficacy of the procedure, solidifying its role as a standard hand surgery.

\section*{Clinical Indications}
De Quervain release is indicated for patients presenting with symptoms and signs consistent with De Quervain tenosynovitis who have failed adequate nonoperative management. Specific clinical indications include:

\begin{itemize}
  \item Persistent, debilitating pain and tenderness over the radial styloid that significantly impairs daily activities, despite a trial of conservative treatment lasting at least six to twelve weeks.
  \item A consistently positive Finkelstein test, where ulnar deviation of the wrist with the thumb flexed into the palm elicits sharp pain over the first dorsal compartment.
  \item Palpable thickening or crepitus over the first dorsal compartment, indicative of chronic inflammation and tendon sheath hypertrophy.
  \item Recurrence of symptoms following one or more corticosteroid injections into the first dorsal compartment, suggesting a persistent mechanical obstruction.
  \item Inability to tolerate or comply with splinting or activity modification due to occupational demands, lifestyle, or patient discomfort.
  \item Suspicion of an anatomical variation, such as a separate subcompartment for the extensor pollicis brevis tendon, which may not respond adequately to injection therapy.
  \item Significant functional limitation in grip strength, pinch strength, and range of motion of the thumb and wrist, impacting quality of life.
  \item Clear clinical diagnosis confirmed by physical examination, ruling out other causes of radial-sided wrist pain such as carpometacarpal arthritis or intersection syndrome.
\end{itemize}

\section*{Clinical Positioning}
De Quervain release is considered a primary surgical treatment for the specific diagnosis of refractory De Quervain tenosynovitis. However, its clinical positioning within the overall treatment algorithm is typically secondary to a thorough trial of nonoperative management. Given that a significant proportion of patients (often 50-70\%) respond favorably to conservative measures such as rest, activity modification, thumb spica splinting, and corticosteroid injections, surgery is not usually the first-line approach.

Instead, it serves as the definitive intervention when conservative strategies have been exhausted or proven ineffective in providing lasting relief. In this context, it occupies a second-line role, becoming the primary definitive treatment once the decision for surgical intervention is made. For a small subset of patients presenting with exceptionally severe, acute symptoms, or those with clear anatomical variations (e.g., a palpable accessory compartment) that are unlikely to respond to injections, surgery might be considered earlier. It is crucial to emphasize that De Quervain release is highly specific; it is not a general solution for radial wrist pain, and careful differential diagnosis is paramount before proceeding with the operation to ensure appropriate patient selection and optimal outcomes.

\section*{Surgical Technique \& Procedure}
The De Quervain release procedure is typically performed on an outpatient basis and can be done under various anesthetic modalities.

\begin{itemize}
  \item \textbf{Anesthesia and Positioning:} The procedure is most commonly performed under local anesthesia with or without intravenous sedation, allowing for patient comfort while maintaining cooperation for nerve monitoring if desired. Regional anesthesia (e.g., axillary block) or general anesthesia may also be used. The patient is positioned supine with the affected arm abducted and externally rotated, resting on a specialized hand table. A tourniquet is applied to the upper arm or forearm and inflated to provide a bloodless surgical field, enhancing visualization of delicate structures.
  \item \textbf{Incision:} A small, typically 2-3 cm, transverse or longitudinal incision is made directly over the first dorsal compartment at the level of the radial styloid. A transverse incision is often preferred for cosmetic reasons, while a longitudinal incision may offer slightly better exposure in complex cases.
  \item \textbf{Dissection and Nerve Protection:} Meticulous dissection proceeds through the subcutaneous tissue. Extreme care is taken to identify and protect the branches of the superficial radial sensory nerve (SRSN), which are highly variable in their course but typically lie superficial and radial to the compartment. These branches are gently retracted to prevent injury.
  \item \textbf{Compartment Exposure:} The thickened extensor retinaculum forming the roof of the first dorsal compartment is identified. It appears as a glistening, fibrous band overlying the APL and EPB tendons.
  \item \textbf{Retinacular Division:} The retinaculum is carefully incised longitudinally along its full length, typically starting proximally and extending distally, ensuring complete release. The incision is made directly over the tendons, taking care not to injure the underlying synovial sheath or tendons themselves. The surgeon confirms that both the APL and EPB tendons are fully decompressed.
  \item \textbf{Accessory Compartment Check:} The compartment is thoroughly inspected for the presence of a separate subcompartment for the EPB tendon or accessory slips of the APL. If present, these must also be released to prevent persistent symptoms. The tendons are gently palpated and moved to ensure smooth, unimpeded gliding.
  \item \textbf{Synovectomy (Optional):} If significant synovial inflammation or hypertrophy is present, a limited synovectomy may be performed to remove the inflamed tissue, although this is not always necessary for successful decompression.
  \item \textbf{Wound Closure:} The tourniquet is deflated, and hemostasis is achieved. The skin is closed with fine sutures or sterile adhesive strips. A soft, bulky dressing is applied, often with a light compression bandage, allowing for immediate gentle range of motion.
\end{itemize}

\section*{Preoperative Workup \& Preparation}
Preoperative workup for De Quervain release is generally straightforward, as the diagnosis is primarily clinical. A thorough medical history and physical examination are paramount to confirm the diagnosis and rule out other causes of radial wrist pain, such as carpometacarpal arthritis, intersection syndrome, or radial nerve entrapment. The surgeon will review the patient's history of conservative treatments, including the duration and efficacy of splinting, activity modification, and corticosteroid injections. If a recent injection has been administered, surgery is typically delayed for several weeks (e.g., 4-6 weeks) to minimize the theoretical risk of infection and allow any residual corticosteroid effects to subside.

\begin{itemize}
  \item \textbf{Medical Evaluation:} A routine preoperative medical evaluation is performed, tailored to the patient's age and comorbidities. This typically includes a review of systems, vital signs, and an assessment of cardiac and pulmonary status.
  \item \textbf{Laboratory Tests:} Basic laboratory tests, such as a complete blood count and coagulation profile, may be ordered, especially if the patient is on anticoagulant medications or has a history of bleeding disorders.
  \item \textbf{Imaging:} Imaging studies are rarely required for diagnosis but may be considered if there is diagnostic uncertainty (e.g., X-rays to rule out fracture or arthritis, ultrasound to confirm tenosynovitis or identify anatomical variations).
  \item \textbf{Medication Adjustment:} Patients on antiplatelet agents (e.g., aspirin, clopidogrel) or anticoagulants (e.g., warfarin, DOACs) will have their medication regimen reviewed. A plan for temporary cessation or bridging therapy will be made in consultation with the prescribing physician, balancing the risk of bleeding against the risk of thrombosis.
  \item \textbf{NPO Status:} Fasting instructions (NPO - nil per os) are provided based on the type of anesthesia planned. For local anesthesia, minimal fasting may be required, while general or regional anesthesia typically necessitates an overnight fast.
  \item \textbf{Prophylactic Antibiotics:} Prophylactic intravenous antibiotics are generally not routinely administered for this clean, superficial procedure unless specific risk factors for infection are present.
  \item \textbf{Informed Consent:} Detailed informed consent is obtained, discussing the procedure, expected outcomes, potential complications (especially superficial radial nerve injury), and the possibility of incomplete relief or recurrence.
  \item \textbf{Site Marking:} On the day of surgery, the operative site is clearly marked by the surgeon, and a final "time-out" procedure is performed to confirm patient identity, procedure, and operative site.
\end{itemize}

\section*{Contraindications \& Precautions}
Careful consideration of contraindications and adherence to precautions are essential for safe and effective De Quervain release.

\textbf{Absolute Contraindications:}
\begin{itemize}
  \item Active infection (e.g., cellulitis, abscess) over the proposed operative site, which must be treated and resolved prior to surgery.
  \item Uncontrolled coagulopathy or bleeding disorder that cannot be safely managed preoperatively.
  \item Uncertain diagnosis where other causes of radial wrist pain have not been adequately excluded.
  \item Patient refusal or inability to provide informed consent.
\end{itemize}

\textbf{Relative Contraindications:}
\begin{itemize}
  \item Very mild symptoms that have not undergone an adequate trial of conservative therapy, as many patients respond to nonoperative measures.
  \item Pregnancy or breastfeeding, where an elective procedure can often be safely deferred until after delivery or cessation of breastfeeding.
  \item Severe medical comorbidities that significantly increase the risks of anesthesia and surgery, making the elective procedure unsafe.
  \item Unrealistic patient expectations regarding the outcome or recovery timeline.
\end{itemize}

\textbf{Precautions:}
\begin{itemize}
  \item Meticulous identification and protection of the superficial radial sensory nerve (SRSN) branches are paramount to prevent iatrogenic injury, which can lead to numbness, dysesthesia, or painful neuroma.
  \item Avoidance of excessive retraction during dissection to minimize trauma to surrounding soft tissues and nerve branches.
  \item Thorough inspection for accessory compartments or anomalous tendon slips, particularly for the extensor pollicis brevis, to ensure complete decompression and prevent persistent symptoms.
  \item If the retinaculum is found to be excessively lax after release, or if the tendons appear prone to subluxation, a partial reconstruction or repair of the retinaculum may be considered to prevent postoperative tendon instability, though this is rare.
  \item Careful hemostasis to prevent hematoma formation, which can cause pain and increase the risk of infection.
\end{itemize}

\section*{Complications \& Risks}
While De Quervain release is generally a safe and effective procedure, like all surgeries, it carries potential complications and risks. These can be broadly categorized into general surgical risks and procedure-specific risks.

\textbf{General Surgical Complications:}
\begin{itemize}
  \item \textbf{Bleeding and Hematoma:} Minor bleeding is common, but significant hematoma formation can cause pain, swelling, and increase the risk of infection.
  \item \textbf{Infection:} Superficial wound infection is a rare but possible complication, typically managed with antibiotics. Deep infection is exceedingly rare.
  \item \textbf{Anesthesia Risks:} Risks associated with anesthesia (local, regional, or general) include allergic reactions, respiratory or cardiac complications, and nausea/vomiting.
  \item \textbf{Pain:} Postoperative pain is expected but usually well-controlled with analgesics. Chronic pain can occur due to various factors.
  \item \textbf{Scarring:} All incisions result in a scar. Hypertrophic or keloid scarring can occur, especially in susceptible individuals, leading to cosmetic concerns or tenderness.
\end{itemize}

\textbf{Procedure-Specific Complications:}
\begin{itemize}
  \item \textbf{Superficial Radial Sensory Nerve (SRSN) Injury:} This is the most common and significant complication, potentially leading to numbness, tingling (dysesthesia), or a painful neuroma on the dorsum of the thumb and radial hand.
  \item \textbf{Incomplete Release:} Failure to fully divide the retinaculum or missing an accessory compartment (e.g., a separate compartment for the EPB tendon) can result in persistent or recurrent symptoms.
  \item \textbf{Tendon Subluxation or Bowstringing:} If too much of the retinaculum is excised, the released tendons (APL and EPB) may subluxate or "bowstring" volarly, causing snapping, discomfort, or instability during wrist and thumb motion.
  \item \textbf{Recurrence of Tenosynovitis:} Although rare after complete release, inflammation and pain can recur, sometimes due to scar tissue formation around the tendons.
  \item \textbf{Stiffness:} Postoperative stiffness of the thumb and wrist joints can occur, particularly if early mobilization is not encouraged or if significant scar tissue forms.
  \item \textbf{Adhesions:} Scar tissue can form between the released tendons and surrounding tissues, limiting tendon gliding and causing pain.
  \item \textbf{Complex Regional Pain Syndrome (CRPS):} A rare but severe complication characterized by chronic pain, swelling, stiffness, and skin changes, disproportionate to the initial injury or surgery.
  \item \textbf{Radial Artery Injury:} Though less common due to its deeper position, injury to the radial artery is a potential risk during deep dissection.
\end{itemize}

\section*{Postoperative Recovery Timeline}
The postoperative recovery from De Quervain release is typically rapid, with most patients experiencing significant improvement within a few weeks.

Immediately after surgery, in the Post-Anesthesia Care Unit (PACU), the patient is monitored for recovery from anesthesia, pain control, and any immediate complications. The hand is usually elevated to minimize swelling, and a light, bulky dressing is applied. Most patients are discharged home the same day. During the first 24-48 hours, pain is managed with oral analgesics, typically acetaminophen and NSAIDs, with stronger medications reserved for breakthrough pain. Patients are encouraged to keep the hand elevated and the dressing clean and dry. Gentle active range of motion of the thumb and wrist is often initiated within the first few days to prevent stiffness and adhesions, as tolerated.

Over the first 1-2 weeks, the initial dressing is usually removed, and the wound is inspected. If non-dissolving sutures were used, they are typically removed at this follow-up visit. Patients are advised to continue gentle exercises and gradually increase activity. Light daily activities, such as dressing and eating, can usually be resumed. Return to desk work or light duties is often possible within 1-2 weeks. Swelling and tenderness around the incision site gradually subside. Long-term recovery typically spans 4-6 weeks, during which time patients progressively increase their activity levels, aiming for full, unrestricted use of the hand and wrist. Most individuals achieve complete resolution of pain and return to all pre-injury activities, including sports and manual labor, without limitation.

\section*{Surgical Findings \& Pathology}
During De Quervain release, the surgeon's intraoperative findings are crucial for confirming the diagnosis and guiding the extent of the release. The most common finding is a visibly thickened and fibrotic extensor retinaculum forming the roof of the first dorsal compartment. This retinaculum may appear pearly white and taut, constricting the underlying tendons. Upon incision of the retinaculum, the abductor pollicis longus (APL) and extensor pollicis brevis (EPB) tendons are exposed. These tendons may show signs of chronic inflammation, such as hyperemia, fraying, or nodularity. Often, the synovial sheath surrounding the tendons is also inflamed and hypertrophied, appearing edematous and sometimes boggy. A critical finding to identify is the presence of an accessory compartment, particularly a separate subcompartment for the EPB tendon, which is found in a significant percentage of patients and must also be released to ensure complete decompression.

If any tissue is resected, such as a portion of the thickened retinaculum or inflamed synovium, it is sent for histopathological examination. The pathology report typically confirms the clinical diagnosis by demonstrating chronic inflammation, fibrosis, and myxoid degeneration within the retinacular tissue. The synovial lining may show hyperplasia and infiltration by chronic inflammatory cells, consistent with tenosynovitis. These findings correlate directly with the mechanical constriction and inflammatory process observed clinically and intraoperatively, providing histological confirmation of the disease.

\section*{Expected Postoperative Course}
A normal and successful postoperative course following De Quervain release is characterized by a progressive and sustained reduction in pain and tenderness along the radial side of the wrist. Patients should experience a gradual return to a full and comfortable range of motion of the thumb and wrist, free from the mechanical catching or friction that was present preoperatively. The Finkelstein test, which was previously positive and painful, should become negative or significantly less painful. Healing of the surgical incision is typically uncomplicated, resulting in a small, well-healed scar that fades over time. Swelling and bruising around the operative site are common in the initial days but should steadily resolve. The tendons should glide smoothly within the decompressed compartment, and the patient should regain normal grip and pinch strength without residual snapping or subluxation. The ultimate goal is for the patient to return to all prior activities, including work, hobbies, and sports, without limitation or recurrence of symptoms, indicating a complete and effective resolution of the tenosynovitis.

\section*{Postoperative Care \& Follow-up}
Postoperative care for De Quervain release focuses on wound healing, pain management, and early mobilization to prevent stiffness.

\begin{itemize}
  \item \textbf{Wound Care:} The initial dressing is typically kept clean and dry for 2-3 days. After this, a smaller dressing or bandage may be applied, and gentle washing of the area can begin.
  \item \textbf{Suture Removal:} If non-dissolving sutures were used, they are usually removed at the first follow-up visit, typically 10-14 days post-surgery.
  \item \textbf{Activity Restrictions:} Patients are encouraged to begin gentle active range of motion exercises for the thumb and wrist almost immediately to prevent adhesions. Heavy lifting, forceful gripping, and repetitive thumb movements are restricted for 4-6 weeks.
  \item \textbf{Pain Management:} Oral analgesics (e.g., NSAIDs, acetaminophen) are usually sufficient for pain control. Ice application and elevation can also help reduce swelling and discomfort.
  \item \textbf{Physical Therapy/Occupational Therapy:} While not always required, a formal course of hand therapy may be recommended for patients experiencing persistent stiffness, weakness, or those with complex anatomical variations. More commonly, a home exercise program is provided.
  \item \textbf{Follow-up Visits:} The first follow-up is typically at 1-2 weeks for wound check and suture removal. Further visits are scheduled as needed, depending on the patient's progress and any developing concerns.
\end{itemize}
Patients are advised to contact their surgeon immediately if they experience signs of infection (increasing pain, fever, spreading redness, pus-like discharge), new or worsening numbness, persistent severe pain, or any signs of tendon subluxation or snapping.

\section*{Epidemiology}
De Quervain tenosynovitis is a relatively common musculoskeletal condition, affecting an estimated 1-2\% of the general population at some point in their lives, with an annual incidence reported to be around 0.5 to 1.3 new cases per 1,000 persons. It exhibits a distinct predilection for women, with a female-to-male ratio ranging from 4:1 to 10:1, and is most frequently observed in individuals between the ages of 30 and 50 years. Certain activities involving repetitive or forceful gripping, pinching, and wrist deviation, such as those found in specific occupations (e.g., assembly line workers, carpenters), hobbies (e.g., knitting, gardening), and particularly in new mothers or caregivers of infants (due to repetitive lifting and holding), are recognized risk factors. While not a life-threatening condition, the chronic pain and functional limitations associated with persistent De Quervain tenosynovitis can significantly impact an individual's quality of life, work productivity, and ability to perform daily activities, making surgical release a commonly performed and important hand operation. Mortality is virtually non-existent, and morbidity is primarily related to the pain and functional impairment of the condition itself or potential surgical complications.

\section*{Outcomes \& Prognosis}
Surgical release of the first dorsal compartment for De Quervain tenosynovitis boasts excellent outcomes and a highly favorable prognosis for properly selected patients. Reported success rates for achieving lasting pain relief and functional improvement range from 80\% to over 95\%. The vast majority of patients express high satisfaction with the surgical results, experiencing complete resolution of their symptoms and a return to full, unrestricted activity within a few weeks to months. Functional outcomes, including grip and pinch strength, typically normalize.

The risk of recurrence after a meticulously performed and complete surgical release is low. When symptoms persist or recur, they are most commonly attributable to an unrecognized accessory compartment (especially for the EPB tendon) that was not released, iatrogenic injury to the superficial radial sensory nerve, or, less commonly, tendon subluxation or significant scar tissue formation. Prognostic factors favoring a good outcome include an accurate preoperative diagnosis, a clear failure of conservative management, complete surgical decompression of all involved tendon slips and compartments, and careful protection of the superficial radial nerve. Patients with long-standing, severe symptoms or those with pre-existing nerve irritation may experience a slightly longer recovery period or less complete resolution of all symptoms, but overall, the procedure offers a highly reliable and effective solution for this debilitating condition.

\section*{References}
\begin{enumerate}
  \item MedlinePlus. De Quervain tenosynovitis surgery. U.S. National Library of Medicine; 2025. Available from: \url{https://medlineplus.gov/ency/article/007268.htm}
  \item Green's Operative Hand Surgery. 8th ed. Philadelphia, PA: Elsevier; 2024. Chapter 46: Tendon Disorders of the Wrist and Hand.
  \item American Academy of Orthopaedic Surgeons (AAOS). De Quervain's Tenosynovitis. OrthoInfo; 2023. Available from: \url{https://orthoinfo.aaos.org/en/diseases--conditions/de-quervains-tenosynovitis/}
  \item Sabiston Textbook of Surgery: The Biological Basis of Modern Surgical Practice. 22nd ed. Philadelphia, PA: Elsevier; 2022. Chapter 67: Hand Surgery.
\end{enumerate}

\clearpage